% !TeX root = ../main.tex
\section{支撑材料与复现说明}

\subsection{问题一复现环境与文件对应}

问题一计算在 Windows 环境下使用 Python 3.12.10 完成，绘图使用 Matplotlib 3.10.9；概率计算为确定性精确枚举，不使用随机数。表~\ref{tab:q1-supporting-files}给出支撑材料与论文结果的对应关系。

\begin{table}[H]
  \centering
  \caption{问题一支撑材料及用途}
  \label{tab:q1-supporting-files}
  \begin{tabularx}{\textwidth}{>{\raggedright\arraybackslash}p{0.42\textwidth}X}
    \toprule
    文件 & 用途 \\
    \midrule
    \path{code/problem1.py} & 二项与超几何概率、临界值搜索、观测判定、敏感性分析和图表生成 \\
    \path{code/test_problem1.py} & 概率归一化、边界值、三态判定和有限总体退化性质测试 \\
    \path{results/q1_summary.json} & 正式参数、最小样本量、临界概率和一致性校核汇总 \\
    \path{results/q1_thresholds.csv} & 样本量 $1$ 至 $200$ 的接收与拒收阈值 \\
    \path{results/q1_sensitivity.csv} & 置信水平敏感性原始结果 \\
    \path{results/q1_population_sensitivity.csv} & 有限总体量检验原始结果 \\
    \bottomrule
  \end{tabularx}
\end{table}

\subsection{问题一核心程序}

下列代码分别给出二项与超几何概率计算，以及固定样本临界值、最小样本量和三态判定的核心实现。完整程序与全部结果文件随支撑材料一并提交。

\lstinputlisting[
  language=Python,
  firstline=119,
  lastline=168,
  caption={问题一的二项与超几何概率计算},
  label={lst:q1-probabilities}
]{../code/problem1.py}

\lstinputlisting[
  language=Python,
  firstline=237,
  lastline=450,
  caption={问题一的临界值搜索与观测判定},
  label={lst:q1-thresholds}
]{../code/problem1.py}

\section{问题二复现说明}

问题二程序从初始联合信念出发生成可达状态，执行值迭代并输出状态策略、动作价值、收敛历史和六种情形汇总；绘图程序在此基础上重算参数扰动与信念离散精度，生成正文三幅 Figure Pack。两段程序均为确定性计算，不使用随机种子。

\begin{table}[H]
  \centering
  \caption{问题二支撑材料及用途}
  \label{tab:q2-supporting-files}
  \begin{tabularx}{\textwidth}{>{\raggedright\arraybackslash}p{0.44\textwidth}X}
    \toprule
    文件 & 用途 \\
    \midrule
    \path{programs/q2_decision_model.py} & 信念状态、贝叶斯转移、Bellman 值迭代及结果导出 \\
    \path{code/problem2_figures.py} & 策略路径、收敛、离散精度和参数敏感性 Figure Pack \\
    \path{code/test_problem2.py} & 转移概率、Bellman 残差、正式成本与受限策略比较测试 \\
    \path{programs/results/q2_summary.json} & 求解器设置、六种情形结果及最大残差汇总 \\
    \path{programs/results/q2_state_policy.csv} & 所有可达信念状态的最优动作 \\
    \path{programs/results/q2_policy_results.csv} & 各状态下全部可行动作的 $Q$ 值 \\
    \path{figures/q2/data/q2_sensitivity.csv} & 三类参数扰动的利润、策略切换与残差 \\
    \path{figures/q2/data/q2_rounding_validation.csv} & 四至七位信念合并精度的成本比较 \\
    \bottomrule
  \end{tabularx}
\end{table}

下列程序段给出动作价值、Bellman 最优动作选择与值迭代的核心实现；完整程序及全部结果随支撑材料提交。

\lstinputlisting[
  language=Python,
  firstline=373,
  lastline=455,
  caption={问题二的 Bellman 动作选择与值迭代},
  label={lst:q2-model}
]{../programs/q2_decision_model.py}

\section{问题三复现说明}

问题三程序按 16 位策略生成可达状态并求解期望成本线性方程组，再用遗传算法搜索，以全枚举全部 $65536$ 个策略作为全局最优基准；绘图程序重跑 20 次 GA 记录逐代收敛历史，并生成利润分布、前 10 名和参数敏感性 Figure Pack，示意脚本按同一布局渲染生产结构图。GA 使用固定种子，全部结果可确定性复现。

\begin{table}[H]
  \centering
  \caption{问题三支撑材料及用途}
  \label{tab:q3-supporting-files}
  \begin{tabularx}{\textwidth}{>{\raggedright\arraybackslash}p{0.44\textwidth}X}
    \toprule
    文件 & 用途 \\
    \midrule
    \path{programs/q3_decision_model.py} & 固定策略期望递推、遗传算法、全枚举与 Monte Carlo 复核 \\
    \path{code/problem3_figures.py} & GA 收敛、利润分布、前 10 与敏感性 Figure Pack \\
    \path{code/problem3_schematic.py} & 生产结构机制图的 PDF/SVG 渲染 \\
    \path{programs/results/q3_summary.txt} & 最优策略、期望成本与利润、GA 命中率、Monte Carlo 汇总 \\
    \path{programs/results/q3_top10_strategies.csv} & 全枚举排名前 10 的策略 \\
    \path{programs/results/q3_exact_all_strategies.csv} & 全部 $65536$ 个策略的精确评价 \\
    \path{figures/q3/data/q3_ga_history.csv} & 20 次 GA 独立运行的逐代最优利润 \\
    \path{figures/q3/data/q3_sensitivity.csv} & 三类参数扰动的 GA 搜索结果 \\
    \bottomrule
  \end{tabularx}
\end{table}

下列程序段分别给出固定策略期望成本求解与遗传算法主循环的核心实现；完整程序及全部结果随支撑材料提交。

\lstinputlisting[
  language=Python,
  firstline=367,
  lastline=409,
  caption={问题三的固定策略期望成本求解},
  label={lst:q3-evaluate}
]{../programs/q3_decision_model.py}

\lstinputlisting[
  language=Python,
  firstline=448,
  lastline=511,
  caption={问题三的遗传算法主循环},
  label={lst:q3-ga}
]{../programs/q3_decision_model.py}
